\documentclass[11pt]{article}

\usepackage[margin=1in]{geometry}
\usepackage[T1]{fontenc}
\usepackage[utf8]{inputenc}
\usepackage{lmodern}
\usepackage{graphicx}
\usepackage{booktabs}
\usepackage{longtable}
\usepackage{array}
\usepackage{amsmath}
\usepackage{amssymb}
\usepackage{xcolor}
\usepackage{hyperref}
\usepackage{float}

\hypersetup{
    colorlinks=true,
    linkcolor=blue,
    urlcolor=blue,
    citecolor=blue
}

\newcommand{\uc}{\ensuremath{\mu\text{C}/\text{cm}^2}}
\newcommand{\mvcm}{\ensuremath{\text{MV}/\text{cm}}}

\title{Multidomain Preisach Hysteresis Results for Metal / Interlayer / AlScN / Metal Stacks}
\author{Krishna Chemudupati}
\date{\today}

\begin{document}
\maketitle

\begin{abstract}
This report documents the active multidomain Preisach hysteresis workflow in the
present repository and summarizes the main trends that emerge from the generated
figure set. The model computes quasistatic polarization--voltage loops for
metal / interlayer / AlScN / metal stacks by combining a multidomain Preisach
representation of the ferroelectric, one-dimensional electrostatics with finite
metal screening, and a self-consistent bisection solver. The main results are
organized around three questions: how the loop changes with interlayer
thickness, how the choice of interlayer material affects loop retention, and
how sensitive the loop is to electrode boundary conditions. The discussion is
kept close to the active figures and tables so that each modeling conclusion is
connected directly to a generated result.
\end{abstract}

\tableofcontents
\clearpage

\section{Purpose of the model}

The active workflow is designed to answer a specific materials question: how
the ferroelectric hysteresis loop changes when the interlayer thickness,
interlayer material, electrode pair, and AlScN thickness are varied. Throughout
this report the plotted quantity is the net ferroelectric polarization of the
AlScN layer. It is therefore the average polarization of the multidomain state
after the stack electrostatics and the domain switching have been solved
self-consistently. It is not the screening charge, the electric displacement,
or a measured current.

At a single applied voltage $V_a$, the model guesses an internal field acting
on the ferroelectric, determines which domains would switch under that field,
computes the corresponding net polarization $P_{\mathrm{FE}}$, passes that
polarization into the electrostatic stack model, and then recomputes the field
that the stack itself would produce. The guessed field is adjusted until the
field used to switch the domains agrees with the field returned by the
electrostatics. The hysteresis loop is the collection of those self-consistent
solutions along the prescribed voltage sweep.

\section{Baseline active workflow}

The active interlayer-thickness figures use a fixed multidomain setup with
15000 domains, a structural disorder distribution centered at $c/a = 1.54$ with
standard deviation $\sigma_{c/a} = 0.04$, the constitutive maps
\[
P_s(c/a) = 333.33(c/a) - 400 \quad \uc
\]
and
\[
E_c(c/a) = 3.16(c/a) - 1.1 \quad \mvcm,
\]
random initial domain states generated from a fixed seed, a quasistatic
bisection solver, and a polarization-space stopping tolerance of 0.5 \uc. Each
run consists of two voltage cycles, although only the final cycle is plotted.
The interlayer sweep used throughout the active figure set is
\[
0,\ 0.1,\ 0.2,\ 0.5,\ 0.8,\ 1,\ 2,\ 3 \text{ nm}.
\]

The dominant electrostatic ingredients are the top and bottom electrode
screening lengths, the top and bottom work functions, the interlayer thickness
and dielectric constant, the ferroelectric thickness and dielectric constant,
and the instantaneous ferroelectric polarization. In the active code the
screening charge is computed from
\begin{equation}
\sigma_s =
\frac{
\dfrac{P_d t_d}{k_d}
+
\dfrac{P_{\mathrm{FE}} t_f}{k_f}
+
\epsilon_0 V_a
}{
\dfrac{\lambda_t}{k_t}
+
\dfrac{\lambda_b}{k_b}
+
\dfrac{t_d}{k_d}
+
\dfrac{t_f}{k_f}
+
\dfrac{t_i}{k_i}
},
\label{eq:sigma-active-report}
\end{equation}
and the internal ferroelectric field is then written as
\begin{equation}
E_{\mathrm{FE}} =
\frac{\sigma_s - P_{\mathrm{FE}}}{k_f \epsilon_0}
+
\frac{\Delta V_{\mathrm{bi,FE}}}{t_f}.
\label{eq:efe-active-report}
\end{equation}
These two relations are enough to explain the trends seen later in the figures.
A thicker interlayer weakens the loop because it both reduces how effectively
the applied voltage reaches the ferroelectric and increases the electrostatic
penalty associated with polarization. A higher-$k$ interlayer weakens the loop
less strongly because it is a better dielectric match to the ferroelectric.
Poorer metal screening reduces the stable polarization on both sides of the
loop, while work-function asymmetry introduces a horizontal preference through
the built-in term.

\section{Al$_2$O$_3$ interlayer sweeps}

To keep the sweeps informative rather than trivially saturated, the active
figure set uses a different voltage window for each AlScN thickness. The
45\,nm case is run at $\pm 20$ V, the 20\,nm case at $\pm 12$ V, and the
10\,nm case at $\pm 6.5$ V. The resulting Al$_2$O$_3$ interlayer sweeps are
shown in Figures~\ref{fig:al2o3-45}, \ref{fig:al2o3-20}, and
\ref{fig:al2o3-10}.

\begin{figure}[p]
\centering
\includegraphics[width=0.82\linewidth]{45nm_figures/Ti_Al2O3_AlScN_Ti_il_hysteresis.png}
\caption{Ti/Al$_2$O$_3$/AlScN/Ti IL-thickness sweep for 45\,nm AlScN at $\pm 20$ V.}
\label{fig:al2o3-45}
\end{figure}

\begin{figure}[p]
\centering
\includegraphics[width=0.82\linewidth]{20nm_figures/Ti_Al2O3_AlScN_Ti_il_hysteresis_12V.png}
\caption{Ti/Al$_2$O$_3$/AlScN/Ti IL-thickness sweep for 20\,nm AlScN at $\pm 12$ V.}
\label{fig:al2o3-20}
\end{figure}

\begin{figure}[p]
\centering
\includegraphics[width=0.82\linewidth]{10nm_figures/Ti_Al2O3_AlScN_Ti_il_hysteresis_6p5V.png}
\caption{Ti/Al$_2$O$_3$/AlScN/Ti IL-thickness sweep for 10\,nm AlScN at $\pm 6.5$ V.}
\label{fig:al2o3-10}
\end{figure}

\clearpage

Table~\ref{tab:al2o3-pmax} summarizes the positive-branch amplitudes extracted
from these loops.

\begin{table}[H]
\centering
\caption{Positive-branch $P_{\max}$ values for Ti/Al$_2$O$_3$/AlScN/Ti loops (\uc).}
\label{tab:al2o3-pmax}
\small
\begin{tabular}{@{}lrrrrrrrr@{}}
\toprule
Case & 0 nm & 0.1 nm & 0.2 nm & 0.5 nm & 0.8 nm & 1 nm & 2 nm & 3 nm \\
\midrule
45 nm, $\pm 20$ V & 112.835 & 95.340 & 72.146 & 38.986 & 25.992 & 21.203 & 9.600 & 5.380 \\
20 nm, $\pm 12$ V & 112.977 & 113.167 & 106.027 & 56.070 & 36.548 & 29.397 & 13.630 & 7.830 \\
10 nm, $\pm 6.5$ V & 112.895 & 101.079 & 68.863 & 33.053 & 20.600 & 16.104 & 6.498 & 3.003 \\
\bottomrule
\end{tabular}
\end{table}

Several points are worth noting. At 0\,nm IL all three thickness cases are
essentially fully polarized, with amplitudes very close to $\pm 113$ \uc. Once
the IL is introduced, the loops shrink rapidly in all three cases. The 45\,nm
case shows the strongest suppression at large IL for the chosen voltage window
and falls to 5.380 \uc\ at 3\,nm IL. The 20\,nm case remains the most robust of
the three under its chosen drive and still reaches 7.830 \uc\ at 3\,nm IL. The
10\,nm case at $\pm 6.5$ V is more clearly voltage-limited and falls to
3.003 \uc\ at 3\,nm IL. The important point is that a thinner ferroelectric is
not simply ``more sensitive'' in a single universal sense. A thinner film is
easier to drive at fixed voltage, but it is also more vulnerable to
depolarization. The apparent thickness trend therefore depends on the voltage
window as well as the geometry itself.

\section{Replacing Al$_2$O$_3$ with HfO$_x$}

The same thickness-specific voltage windows were then applied to
Ti/HfO$_x$/AlScN/Ti. The resulting sweeps are shown in
Figures~\ref{fig:hfox-45}, \ref{fig:hfox-20}, and \ref{fig:hfox-10}.

\begin{figure}[p]
\centering
\includegraphics[width=0.82\linewidth]{45nm_figures/Ti_HfOx_AlScN_Ti_il_hysteresis_20V.png}
\caption{Ti/HfO$_x$/AlScN/Ti IL-thickness sweep for 45\,nm AlScN at $\pm 20$ V.}
\label{fig:hfox-45}
\end{figure}

\begin{figure}[p]
\centering
\includegraphics[width=0.82\linewidth]{20nm_figures/Ti_HfOx_AlScN_Ti_il_hysteresis_12V.png}
\caption{Ti/HfO$_x$/AlScN/Ti IL-thickness sweep for 20\,nm AlScN at $\pm 12$ V.}
\label{fig:hfox-20}
\end{figure}

\begin{figure}[p]
\centering
\includegraphics[width=0.82\linewidth]{10nm_figures/Ti_HfOx_AlScN_Ti_il_hysteresis_6p5V.png}
\caption{Ti/HfO$_x$/AlScN/Ti IL-thickness sweep for 10\,nm AlScN at $\pm 6.5$ V.}
\label{fig:hfox-10}
\end{figure}

\clearpage

\begin{table}[H]
\centering
\caption{Positive-branch $P_{\max}$ values for Ti/HfO$_x$/AlScN/Ti loops (\uc).}
\label{tab:hfox-pmax}
\small
\begin{tabular}{@{}lrrrrrrrr@{}}
\toprule
Case & 0 nm & 0.1 nm & 0.2 nm & 0.5 nm & 0.8 nm & 1 nm & 2 nm & 3 nm \\
\midrule
45 nm, $\pm 20$ V & 112.835 & 106.786 & 92.224 & 59.221 & 42.542 & 35.610 & 18.843 & 11.999 \\
20 nm, $\pm 12$ V & 112.977 & 113.167 & 113.167 & 87.398 & 61.313 & 50.687 & 26.230 & 16.533 \\
10 nm, $\pm 6.5$ V & 112.895 & 113.188 & 95.970 & 54.278 & 36.693 & 29.793 & 14.119 & 8.426 \\
\bottomrule
\end{tabular}
\end{table}

The dominant result is that, for the same geometry and voltage window, HfO$_x$
preserves much larger loops than Al$_2$O$_3$ at finite IL thickness. This
follows directly from the active dielectric constants:
$k_{\mathrm{AlScN}} = 16$, $k_{\mathrm{Al_2O_3}} = 9.3$, and
$k_{\mathrm{HfO_x}} = 16.64$. HfO$_x$ is therefore a much better electrostatic
match to AlScN than Al$_2$O$_3$. In the present model this means that less
voltage is lost across the IL and the depolarization penalty is weaker. The
effect is substantial. For the 45\,nm case at 3\,nm IL, the HfO$_x$ loop
reaches 11.999 \uc\ while the Al$_2$O$_3$ loop reaches only 5.380 \uc. For the
20\,nm case at the same IL thickness, HfO$_x$ gives 16.533 \uc\ while
Al$_2$O$_3$ gives 7.830 \uc. For the 10\,nm case the corresponding values are
8.426 \uc\ and 3.003 \uc. In other words, switching from Al$_2$O$_3$ to HfO$_x$
roughly doubles the large-IL loop amplitude in the thicker two cases and
improves it by almost a factor of three in the 10\,nm, lower-voltage case.

\section{Electrode sensitivity at 0\,nm IL}

The repository also includes a 0\,nm IL electrode study, which is useful
because it isolates the electrode boundary condition from interlayer-material
effects. The resulting overlay and summary heatmap are shown in
Figures~\ref{fig:electrode-overlay} and \ref{fig:electrode-heatmap}.

\begin{figure}[p]
\centering
\includegraphics[width=0.94\linewidth]{45nm_figures/all_electrodes_0nm_overlay.png}
\caption{Overlay of all ordered top/bottom electrode combinations at 0\,nm IL. Color denotes the top electrode, while line style and marker denote the bottom electrode.}
\label{fig:electrode-overlay}
\end{figure}

\begin{figure}[p]
\centering
\includegraphics[width=0.75\linewidth]{45nm_figures/all_electrodes_0nm_deltaP_summary.png}
\caption{Mean absolute polarization difference relative to Ti/Ti at 0\,nm IL.}
\label{fig:electrode-heatmap}
\end{figure}

\clearpage

\begin{table}[H]
\centering
\caption{Representative 0\,nm IL electrode effects relative to Ti/Ti.}
\label{tab:electrode-deltas}
\begin{tabular}{@{}lrrrr@{}}
\toprule
Stack & mean $|\Delta P|$ (\uc) & max $|\Delta P|$ (\uc) & $\Delta P(0\text{V},\mathrm{fwd})$ & $\Delta P(0\text{V},\mathrm{rev})$ \\
\midrule
Ti/Al & 0.343 & 2.933 & -0.102 & 0.041 \\
Ti/Pd & 13.014 & 59.715 & 16.930 & -0.041 \\
Pd/Ti & 13.193 & 59.642 & 0.298 & -17.135 \\
Ti/Test & 67.623 & 84.359 & 66.823 & -73.976 \\
Ti/Test2 & 67.365 & 80.180 & 70.329 & -69.851 \\
\bottomrule
\end{tabular}
\end{table}

This comparison separates two physically distinct effects. The first is
work-function asymmetry. Ti/Al remains close to Ti/Ti because the work-function
mismatch is small in the active parameter set. Ti/Pd and Pd/Ti differ much more
strongly, and those two stacks reverse one another in the way expected from a
built-in field that changes sign when the electrode order is reversed. The
second effect is screening-length sensitivity. The most revealing comparison is
Ti/Test2. Test2 was designed to keep the Ti work function while increasing the
screening length, and the fact that Ti/Test2 still differs strongly from Ti/Ti
shows that, in the present hysteresis-only implementation, electrode screening
length can change the loop just as strongly as work-function asymmetry. This is
useful because it shows that the metal boundary condition is not determined by
work function alone; it also depends on how effectively the electrodes
compensate ferroelectric bound charge.

\section{Interpretation of the generated loop shapes}

Some of the generated loops are more boxy or rhomboid than the smooth textbook
shape often associated with ferroelectrics. In the present model this is a
natural consequence of the assumptions rather than a plotting artifact. Each
domain is binary rather than continuously varying, switching is threshold-based
and abrupt, the solver is quasistatic rather than dynamical, and the internal
field can vary almost linearly with applied voltage over part of the sweep.
When a large collection of domains is driven through a fairly ordered set of
thresholds under those conditions, the net polarization can follow straight
segments and produce the boxier loop shape seen in some figures.

\section{Scope and limitations}

The current repository already accounts coherently for several important trends.
It captures the strong suppression caused by increasing interlayer thickness, it
shows why a higher-$k$ interlayer preserves the loop more effectively, it
resolves loop skew introduced by work-function asymmetry, and it demonstrates
that poorer electrode screening can shrink the loop even when the work function
is held fixed. Within the present workflow these trends are emergent outputs of
the coupled equations once the geometry, material set, and sweep protocol have
been fixed.

At the same time, this remains a compact hysteresis model rather than a full
ferroelectric-diode transport simulator. It does not yet include DC current
transport, thermionic emission, tunneling, trap-assisted current mechanisms,
explicit Landau--Khalatnikov dynamics, time-dependent switching kinetics,
spatial domain-wall physics, or finite metal thickness in the electrostatics.

\section{Concluding remarks}

The active multidomain Preisach workflow supports a clear set of conclusions.
The hysteresis curves are not drawn by hand or imposed by post-processing; they
are produced by a coupled memory-plus-electrostatics model. Interlayer
thickness is a dominant control variable because it enters the electrostatic
denominator directly. HfO$_x$ is considerably more favorable than Al$_2$O$_3$
for preserving loop amplitude at finite IL thickness in the present model.
Electrode screening length is a major hysteresis lever alongside work-function
asymmetry. Finally, thickness trends cannot be interpreted independently of the
chosen voltage window, because the ease of driving the ferroelectric and the
depolarization penalty compete throughout the sweep.

\end{document}
